%%%%%%%%%%%%%%%%%%%%%%%%%%%%%%%%%%%%%%%%%%%%%%%%%%%%%%%%%%%%
% 9.7 DELAY DIFFERENTIAL EQUATIONS
% The Composition of Advanced Mathematics
%%%%%%%%%%%%%%%%%%%%%%%%%%%%%%%%%%%%%%%%%%%%%%%%%%%%%%%%%%%%

\section{Delay Differential Equations}
\label{sec:delay-differential-equations}

\prerequisite{
Ordinary differential equations, dynamical systems, stability,
linearization, characteristic equations, Laplace transforms, integral
equations, numerical methods, and basic complex-number methods.
}

\learningobjectives{
By the end of this section, the reader should be able to:
\begin{itemize}
    \item define a delay differential equation;
    \item explain how delay equations differ from ordinary differential
          equations;
    \item distinguish discrete, distributed, constant, and
          state-dependent delays;
    \item understand why delay equations require history functions rather
          than single initial values;
    \item formulate delay initial-value problems;
    \item understand the method of steps;
    \item solve elementary delay equations piecewise;
    \item identify equilibria of autonomous delay equations;
    \item linearize nonlinear delay equations near equilibria;
    \item derive characteristic equations containing exponentials;
    \item understand why delay characteristic equations possess
          infinitely many roots;
    \item analyze stability of scalar linear delay equations;
    \item understand delay-induced instability;
    \item recognize oscillatory instability;
    \item derive imaginary-axis crossing conditions;
    \item understand Hopf bifurcation caused by delay;
    \item understand the Lambert \(W\) representation for simple delay
          equations;
    \item recognize neutral delay equations;
    \item distinguish retarded and neutral equations;
    \item understand distributed-delay models;
    \item connect distributed delays with integral equations;
    \item recognize state-dependent delays;
    \item understand delayed logistic population models;
    \item understand delayed feedback systems;
    \item understand delayed epidemic and physiological models;
    \item formulate systems of DDEs;
    \item understand phase space for delay equations;
    \item recognize that DDEs are effectively infinite-dimensional;
    \item understand numerical approximation of DDEs;
    \item recognize interpolation requirements in numerical DDE solvers;
    \item understand discontinuity propagation in derivatives;
    \item recognize stability restrictions created by delay;
    \item connect DDEs with ODEs, dynamical systems, integral equations,
          control, and mathematical modeling.
\end{itemize}
}

%%%%%%%%%%%%%%%%%%%%%%%%%%%%%%%%%%%%%%%%%%%%%%%%%%%%%%%%%%%%
\subsection{What Is a Delay Differential Equation?}
\label{subsec:what-is-delay-differential-equation}
%%%%%%%%%%%%%%%%%%%%%%%%%%%%%%%%%%%%%%%%%%%%%%%%%%%%%%%%%%%%

A delay differential equation, abbreviated DDE, is a differential
equation in which the rate of change at the present time depends on one
or more past states.

A simple example is

\[
\boxed{
x'(t)
=
-x(t-1).
}
\]

The derivative at time \(t\) depends on the value one time unit earlier.

A more general form is

\[
\boxed{
x'(t)
=
f
\left(
t,
x(t),
x(t-\tau)
\right),
}
\]

where

\[
\boxed{
\tau>0
}
\]

is the delay.

%%%%%%%%%%%%%%%%%%%%%%%%%%%%%%%%%%%%%%%%%%%%%%%%%%%%%%%%%%%%
\subsection{Why Delays Matter}
\label{subsec:why-delays-matter}
%%%%%%%%%%%%%%%%%%%%%%%%%%%%%%%%%%%%%%%%%%%%%%%%%%%%%%%%%%%%

Many real systems do not react instantaneously.

Examples include:

\[
\boxed{
\text{biological maturation},
}
\]

\[
\boxed{
\text{transport time},
}
\]

\[
\boxed{
\text{communication lag},
}
\]

\[
\boxed{
\text{incubation periods},
}
\]

and

\[
\boxed{
\text{delayed feedback}.
}
\]

If the response time is not negligible, an ordinary differential equation
may fail to represent the dynamics accurately.

%%%%%%%%%%%%%%%%%%%%%%%%%%%%%%%%%%%%%%%%%%%%%%%%%%%%%%%%%%%%
\subsection{ODE Versus DDE}
\label{subsec:ode-versus-dde}
%%%%%%%%%%%%%%%%%%%%%%%%%%%%%%%%%%%%%%%%%%%%%%%%%%%%%%%%%%%%

An ODE has present-state dependence such as

\[
\boxed{
x'(t)
=
f(t,x(t)).
}
\]

A DDE may include past-state dependence:

\[
\boxed{
x'(t)
=
f
\left(
t,
x(t),
x(t-\tau)
\right).
}
\]

Thus:

\[
\boxed{
\text{ODE}
\rightarrow
\text{present state},
}
\]

while

\[
\boxed{
\text{DDE}
\rightarrow
\text{present and past states}.
}
\]

%%%%%%%%%%%%%%%%%%%%%%%%%%%%%%%%%%%%%%%%%%%%%%%%%%%%%%%%%%%%
\subsection{Discrete Delay}
\label{subsec:discrete-delay}
%%%%%%%%%%%%%%%%%%%%%%%%%%%%%%%%%%%%%%%%%%%%%%%%%%%%%%%%%%%%

A discrete or point delay uses one or more specific past times.

For example,

\[
\boxed{
x'(t)
=
ax(t)
+
bx(t-\tau).
}
\]

The past contribution is evaluated exactly at

\[
t-\tau.
\]

%%%%%%%%%%%%%%%%%%%%%%%%%%%%%%%%%%%%%%%%%%%%%%%%%%%%%%%%%%%%
\subsection{Multiple Discrete Delays}
\label{subsec:multiple-discrete-delays}
%%%%%%%%%%%%%%%%%%%%%%%%%%%%%%%%%%%%%%%%%%%%%%%%%%%%%%%%%%%%

A system may depend on several past times:

\[
\boxed{
x'(t)
=
ax(t)
+
b_1x(t-\tau_1)
+
b_2x(t-\tau_2).
}
\]

More generally,

\[
\boxed{
x'(t)
=
f
\left(
x(t),
x(t-\tau_1),
\ldots,
x(t-\tau_m)
\right).
}
\]

%%%%%%%%%%%%%%%%%%%%%%%%%%%%%%%%%%%%%%%%%%%%%%%%%%%%%%%%%%%%
\subsection{Constant Delay}
\label{subsec:constant-delay}
%%%%%%%%%%%%%%%%%%%%%%%%%%%%%%%%%%%%%%%%%%%%%%%%%%%%%%%%%%%%

A constant delay satisfies

\[
\boxed{
\tau=\text{constant}.
}
\]

For example,

\[
\boxed{
x'(t)
=
-x(t-2).
}
\]

Constant-delay equations are among the most tractable DDEs.

%%%%%%%%%%%%%%%%%%%%%%%%%%%%%%%%%%%%%%%%%%%%%%%%%%%%%%%%%%%%
\subsection{Time-Dependent Delay}
\label{subsec:time-dependent-delay}
%%%%%%%%%%%%%%%%%%%%%%%%%%%%%%%%%%%%%%%%%%%%%%%%%%%%%%%%%%%%

The delay itself may depend on time:

\[
\boxed{
x'(t)
=
f
\left(
t,
x(t),
x(t-\tau(t))
\right).
}
\]

The quantity

\[
\tau(t)
\]

may vary because of changing transport, communication, or biological
conditions.

%%%%%%%%%%%%%%%%%%%%%%%%%%%%%%%%%%%%%%%%%%%%%%%%%%%%%%%%%%%%
\subsection{State-Dependent Delay}
\label{subsec:state-dependent-delay}
%%%%%%%%%%%%%%%%%%%%%%%%%%%%%%%%%%%%%%%%%%%%%%%%%%%%%%%%%%%%

A state-dependent delay depends on the evolving solution:

\[
\boxed{
x'(t)
=
f
\left(
x(t),
x
\left(
t-\tau(x(t))
\right)
\right).
}
\]

Such equations can be substantially more difficult because the delayed
evaluation time is itself unknown until the solution is known.

%%%%%%%%%%%%%%%%%%%%%%%%%%%%%%%%%%%%%%%%%%%%%%%%%%%%%%%%%%%%
\subsection{Distributed Delay}
\label{subsec:distributed-delay}
%%%%%%%%%%%%%%%%%%%%%%%%%%%%%%%%%%%%%%%%%%%%%%%%%%%%%%%%%%%%

Instead of depending on one past time, a distributed-delay model averages
over a range of past states:

\[
\boxed{
x'(t)
=
f(x(t))
+
\int_0^\tau
K(s)x(t-s)\,ds.
}
\]

The kernel

\[
K(s)
\]

weights the influence of different time lags.

%%%%%%%%%%%%%%%%%%%%%%%%%%%%%%%%%%%%%%%%%%%%%%%%%%%%%%%%%%%%
\subsection{Infinite Distributed Delay}
\label{subsec:infinite-distributed-delay}
%%%%%%%%%%%%%%%%%%%%%%%%%%%%%%%%%%%%%%%%%%%%%%%%%%%%%%%%%%%%

Some models use the entire past:

\[
\boxed{
x'(t)
=
f(x(t))
+
\int_0^\infty
K(s)x(t-s)\,ds.
}
\]

For the integral to be meaningful, the memory kernel usually decays
sufficiently rapidly.

This connects DDEs directly with Volterra integral equations.

%%%%%%%%%%%%%%%%%%%%%%%%%%%%%%%%%%%%%%%%%%%%%%%%%%%%%%%%%%%%
\subsection{Retarded Functional Differential Equations}
\label{subsec:retarded-functional-differential-equations}
%%%%%%%%%%%%%%%%%%%%%%%%%%%%%%%%%%%%%%%%%%%%%%%%%%%%%%%%%%%%

A common abstract form is

\[
\boxed{
x'(t)
=
F(x_t),
}
\]

where \(x_t\) denotes the recent history segment

\[
\boxed{
x_t(\theta)
=
x(t+\theta),
\qquad
-\tau\leq\theta\leq0.
}
\]

The derivative depends on the history but not on delayed derivatives.

Such equations are called retarded functional differential equations.

%%%%%%%%%%%%%%%%%%%%%%%%%%%%%%%%%%%%%%%%%%%%%%%%%%%%%%%%%%%%
\subsection{History Functions}
\label{subsec:dde-history-functions}
%%%%%%%%%%%%%%%%%%%%%%%%%%%%%%%%%%%%%%%%%%%%%%%%%%%%%%%%%%%%

An ODE initial-value problem typically specifies one value:

\[
x(0)=x_0.
\]

A DDE with maximum delay \(\tau\) generally requires an entire history
function:

\[
\boxed{
x(t)
=
\phi(t),
\qquad
-\tau\leq t\leq0.
}
\]

Thus the initial data are themselves a function.

%%%%%%%%%%%%%%%%%%%%%%%%%%%%%%%%%%%%%%%%%%%%%%%%%%%%%%%%%%%%
\subsection{Why One Initial Value Is Not Enough}
\label{subsec:dde-one-value-not-enough}
%%%%%%%%%%%%%%%%%%%%%%%%%%%%%%%%%%%%%%%%%%%%%%%%%%%%%%%%%%%%

Consider

\[
x'(t)
=
x(t-1).
\]

To compute

\[
x'(0.4),
\]

one needs

\[
x(-0.6).
\]

Therefore specifying only

\[
x(0)
\]

does not provide enough information.

One must know the history on

\[
\boxed{
[-1,0].
}
\]

%%%%%%%%%%%%%%%%%%%%%%%%%%%%%%%%%%%%%%%%%%%%%%%%%%%%%%%%%%%%
\subsection{DDE Initial-Value Problem}
\label{subsec:dde-initial-value-problem}
%%%%%%%%%%%%%%%%%%%%%%%%%%%%%%%%%%%%%%%%%%%%%%%%%%%%%%%%%%%%

A basic DDE initial-value problem is

\[
\boxed{
x'(t)
=
f
\left(
t,
x(t),
x(t-\tau)
\right),
\qquad
t>0,
}
\]

with history

\[
\boxed{
x(t)
=
\phi(t),
\qquad
-\tau\leq t\leq0.
}
\]

The function \(\phi\) determines the past state before evolution begins.

%%%%%%%%%%%%%%%%%%%%%%%%%%%%%%%%%%%%%%%%%%%%%%%%%%%%%%%%%%%%
\subsection{Phase Space for Delay Equations}
\label{subsec:dde-phase-space}
%%%%%%%%%%%%%%%%%%%%%%%%%%%%%%%%%%%%%%%%%%%%%%%%%%%%%%%%%%%%

For an ODE in \(\mathbb{R}^n\), the state at one time is a vector.

For a DDE, the state must contain a history segment.

A natural phase space is therefore a function space such as

\[
\boxed{
C
\left(
[-\tau,0],
\mathbb{R}^n
\right).
}
\]

Thus even a scalar DDE is effectively an infinite-dimensional dynamical
system.

%%%%%%%%%%%%%%%%%%%%%%%%%%%%%%%%%%%%%%%%%%%%%%%%%%%%%%%%%%%%
\subsection{Infinite-Dimensional Nature of DDEs}
\label{subsec:dde-infinite-dimensional}
%%%%%%%%%%%%%%%%%%%%%%%%%%%%%%%%%%%%%%%%%%%%%%%%%%%%%%%%%%%%

The present state cannot generally be described by finitely many numbers.

Instead, one needs the entire function

\[
\boxed{
x_t(\theta)
=
x(t+\theta).
}
\]

This explains why DDEs can have much richer behavior than scalar ODEs.

For example, a single scalar DDE can exhibit oscillations and complicated
dynamics that would require higher dimension in ordinary differential
equations.

%%%%%%%%%%%%%%%%%%%%%%%%%%%%%%%%%%%%%%%%%%%%%%%%%%%%%%%%%%%%
\subsection{Method of Steps}
\label{subsec:method-of-steps}
%%%%%%%%%%%%%%%%%%%%%%%%%%%%%%%%%%%%%%%%%%%%%%%%%%%%%%%%%%%%

For constant-delay equations, a basic solution technique is the method of
steps.

Suppose

\[
x'(t)
=
f
\left(
t,
x(t),
x(t-\tau)
\right)
\]

with known history on

\[
[-\tau,0].
\]

Then:

\begin{enumerate}
    \item solve on \(0\leq t\leq\tau\) using the known history;
    \item use that solution as known delayed data on
          \(\tau\leq t\leq2\tau\);
    \item continue interval by interval.
\end{enumerate}

%%%%%%%%%%%%%%%%%%%%%%%%%%%%%%%%%%%%%%%%%%%%%%%%%%%%%%%%%%%%
\subsection{Worked Example: Method of Steps}
\label{subsec:worked-method-of-steps}
%%%%%%%%%%%%%%%%%%%%%%%%%%%%%%%%%%%%%%%%%%%%%%%%%%%%%%%%%%%%

\begin{workedexamplebox}{A Simple Delayed Equation}

Solve

\[
x'(t)=x(t-1),
\qquad
t>0,
\]

with history

\[
x(t)=1,
\qquad
-1\leq t\leq0.
\]

For

\[
0\leq t\leq1,
\]

we have

\[
x(t-1)=1.
\]

Thus,

\[
x'(t)=1.
\]

Using

\[
x(0)=1,
\]

we obtain

\[
\boxed{
x(t)=1+t,
\qquad
0\leq t\leq1.
}
\]

For

\[
1\leq t\leq2,
\]

the delayed value is

\[
x(t-1)
=
1+(t-1)
=
t.
\]

Therefore,

\[
x'(t)=t.
\]

Integrating from \(1\),

\[
x(t)
=
x(1)
+
\int_1^t s\,ds.
\]

Since

\[
x(1)=2,
\]

\[
x(t)
=
2
+
\frac12(t^2-1).
\]

Hence,

\begin{answerbox}
\[
\boxed{
x(t)
=
\frac32
+
\frac12t^2,
\qquad
1\leq t\leq2.
}
\]
\end{answerbox}

The procedure continues step by step.

\end{workedexamplebox}

%%%%%%%%%%%%%%%%%%%%%%%%%%%%%%%%%%%%%%%%%%%%%%%%%%%%%%%%%%%%
\subsection{Piecewise Smoothness}
\label{subsec:dde-piecewise-smoothness}
%%%%%%%%%%%%%%%%%%%%%%%%%%%%%%%%%%%%%%%%%%%%%%%%%%%%%%%%%%%%

Even when the history function is continuous, derivatives of a DDE
solution may develop changes in smoothness at times

\[
\boxed{
\tau,
\ 2\tau,
\ 3\tau,
\ldots
}
\]

because delayed information propagates forward interval by interval.

This is an important issue for numerical solvers.

%%%%%%%%%%%%%%%%%%%%%%%%%%%%%%%%%%%%%%%%%%%%%%%%%%%%%%%%%%%%
\subsection{Compatibility at the Initial Time}
\label{subsec:dde-initial-compatibility}
%%%%%%%%%%%%%%%%%%%%%%%%%%%%%%%%%%%%%%%%%%%%%%%%%%%%%%%%%%%%

Suppose the history is differentiable.

The left derivative at \(t=0\) is

\[
\phi'(0^-).
\]

The DDE determines the right derivative:

\[
\boxed{
x'(0^+)
=
f
\left(
0,
\phi(0),
\phi(-\tau)
\right).
}
\]

If these do not agree, the solution can remain continuous while its
derivative has a jump at \(t=0\).

%%%%%%%%%%%%%%%%%%%%%%%%%%%%%%%%%%%%%%%%%%%%%%%%%%%%%%%%%%%%
\subsection{Existence and Uniqueness}
\label{subsec:dde-existence-uniqueness}
%%%%%%%%%%%%%%%%%%%%%%%%%%%%%%%%%%%%%%%%%%%%%%%%%%%%%%%%%%%%

For a retarded DDE

\[
x'(t)=F(x_t),
\]

existence and uniqueness are obtained under conditions analogous to ODE
theory.

Roughly, continuity provides existence while a suitable Lipschitz
condition in the history variable provides uniqueness.

The precise hypotheses are formulated in function spaces.

%%%%%%%%%%%%%%%%%%%%%%%%%%%%%%%%%%%%%%%%%%%%%%%%%%%%%%%%%%%%
\subsection{Autonomous Delay Equations}
\label{subsec:autonomous-delay-equations}
%%%%%%%%%%%%%%%%%%%%%%%%%%%%%%%%%%%%%%%%%%%%%%%%%%%%%%%%%%%%

An autonomous delay equation has no explicit time dependence.

For example,

\[
\boxed{
x'(t)
=
f
\left(
x(t),
x(t-\tau)
\right).
}
\]

Equilibria are constant solutions.

%%%%%%%%%%%%%%%%%%%%%%%%%%%%%%%%%%%%%%%%%%%%%%%%%%%%%%%%%%%%
\subsection{Equilibria of Delay Equations}
\label{subsec:dde-equilibria}
%%%%%%%%%%%%%%%%%%%%%%%%%%%%%%%%%%%%%%%%%%%%%%%%%%%%%%%%%%%%

Suppose

\[
x(t)=x^*
\]

for all \(t\).

Then

\[
x(t-\tau)=x^*.
\]

Therefore an equilibrium satisfies

\[
\boxed{
f(x^*,x^*)=0.
}
\]

Notice that the delay itself often does not change the location of the
equilibrium, but it can drastically change its stability.

%%%%%%%%%%%%%%%%%%%%%%%%%%%%%%%%%%%%%%%%%%%%%%%%%%%%%%%%%%%%
\subsection{Linear Scalar Delay Equation}
\label{subsec:linear-scalar-delay-equation}
%%%%%%%%%%%%%%%%%%%%%%%%%%%%%%%%%%%%%%%%%%%%%%%%%%%%%%%%%%%%

A basic linear DDE is

\[
\boxed{
x'(t)
=
ax(t)
+
bx(t-\tau).
}
\]

To study stability, try an exponential solution

\[
\boxed{
x(t)=e^{\lambda t}.
}
\]

Then

\[
x'(t)
=
\lambda e^{\lambda t},
\]

and

\[
x(t-\tau)
=
e^{\lambda(t-\tau)}
=
e^{-\lambda\tau}e^{\lambda t}.
\]

%%%%%%%%%%%%%%%%%%%%%%%%%%%%%%%%%%%%%%%%%%%%%%%%%%%%%%%%%%%%
\subsection{Characteristic Equation of a Linear DDE}
\label{subsec:dde-characteristic-equation}
%%%%%%%%%%%%%%%%%%%%%%%%%%%%%%%%%%%%%%%%%%%%%%%%%%%%%%%%%%%%

Substitution into

\[
x'
=
ax+bx(t-\tau)
\]

gives

\[
\lambda e^{\lambda t}
=
ae^{\lambda t}
+
be^{-\lambda\tau}e^{\lambda t}.
\]

Cancel \(e^{\lambda t}\):

\[
\boxed{
\lambda
=
a
+
be^{-\lambda\tau}.
}
\]

Equivalently,

\[
\boxed{
\lambda-a-be^{-\lambda\tau}=0.
}
\]

This is a transcendental characteristic equation.

%%%%%%%%%%%%%%%%%%%%%%%%%%%%%%%%%%%%%%%%%%%%%%%%%%%%%%%%%%%%
\subsection{Why There Are Infinitely Many Characteristic Roots}
\label{subsec:dde-infinitely-many-roots}
%%%%%%%%%%%%%%%%%%%%%%%%%%%%%%%%%%%%%%%%%%%%%%%%%%%%%%%%%%%%

For an ODE, the characteristic equation is usually polynomial.

For a DDE,

\[
e^{-\lambda\tau}
\]

appears.

Because exponential and algebraic terms interact, the equation generally
has infinitely many complex roots.

This is another manifestation of the infinite-dimensional nature of DDEs.

%%%%%%%%%%%%%%%%%%%%%%%%%%%%%%%%%%%%%%%%%%%%%%%%%%%%%%%%%%%%
\subsection{Stability from Characteristic Roots}
\label{subsec:dde-characteristic-stability}
%%%%%%%%%%%%%%%%%%%%%%%%%%%%%%%%%%%%%%%%%%%%%%%%%%%%%%%%%%%%

For a linear retarded DDE, exponential stability is governed by the real
parts of the characteristic roots.

A standard criterion is:

\[
\boxed{
\operatorname{Re}(\lambda)<0
\quad
\text{for every characteristic root}
}
\]

for asymptotic or exponential stability under the usual hypotheses.

If any root satisfies

\[
\boxed{
\operatorname{Re}(\lambda)>0,
}
\]

the equilibrium is unstable.

%%%%%%%%%%%%%%%%%%%%%%%%%%%%%%%%%%%%%%%%%%%%%%%%%%%%%%%%%%%%
\subsection{No-Delay Limit}
\label{subsec:dde-no-delay-limit}
%%%%%%%%%%%%%%%%%%%%%%%%%%%%%%%%%%%%%%%%%%%%%%%%%%%%%%%%%%%%

If

\[
\tau=0,
\]

then

\[
x(t-\tau)=x(t).
\]

Thus

\[
x'
=
ax+bx
=
(a+b)x.
\]

The ODE stability condition is

\[
\boxed{
a+b<0.
}
\]

Increasing \(\tau\) can destabilize an equilibrium that is stable when
\(\tau=0\).

%%%%%%%%%%%%%%%%%%%%%%%%%%%%%%%%%%%%%%%%%%%%%%%%%%%%%%%%%%%%
\subsection{Pure Delayed Negative Feedback}
\label{subsec:pure-delayed-negative-feedback}
%%%%%%%%%%%%%%%%%%%%%%%%%%%%%%%%%%%%%%%%%%%%%%%%%%%%%%%%%%%%

Consider

\[
\boxed{
x'(t)
=
-kx(t-\tau),
\qquad
k>0.
}
\]

The characteristic equation is

\[
\boxed{
\lambda
+
ke^{-\lambda\tau}
=
0.
}
\]

This simple equation illustrates delay-induced oscillation.

%%%%%%%%%%%%%%%%%%%%%%%%%%%%%%%%%%%%%%%%%%%%%%%%%%%%%%%%%%%%
\subsection{Imaginary-Axis Crossing}
\label{subsec:dde-imaginary-axis-crossing}
%%%%%%%%%%%%%%%%%%%%%%%%%%%%%%%%%%%%%%%%%%%%%%%%%%%%%%%%%%%%

Stability can change when characteristic roots cross the imaginary axis.

Set

\[
\boxed{
\lambda=i\omega.
}
\]

For

\[
\lambda
+
ke^{-\lambda\tau}
=
0,
\]

we obtain

\[
i\omega
+
k
e^{-i\omega\tau}
=
0.
\]

Using

\[
e^{-i\omega\tau}
=
\cos(\omega\tau)
-
i\sin(\omega\tau),
\]

we get

\[
i\omega
+
k\cos(\omega\tau)
-
ik\sin(\omega\tau)
=
0.
\]

%%%%%%%%%%%%%%%%%%%%%%%%%%%%%%%%%%%%%%%%%%%%%%%%%%%%%%%%%%%%
\subsection{Separating Real and Imaginary Parts}
\label{subsec:dde-real-imaginary-parts}
%%%%%%%%%%%%%%%%%%%%%%%%%%%%%%%%%%%%%%%%%%%%%%%%%%%%%%%%%%%%

The real part gives

\[
\boxed{
k\cos(\omega\tau)=0.
}
\]

Thus,

\[
\boxed{
\cos(\omega\tau)=0.
}
\]

The imaginary part gives

\[
\boxed{
\omega
-
k\sin(\omega\tau)
=
0.
}
\]

At the first positive crossing,

\[
\omega\tau
=
\frac{\pi}{2}.
\]

Then

\[
\sin(\omega\tau)=1,
\]

so

\[
\boxed{
\omega=k.
}
\]

Therefore,

\[
\boxed{
k\tau
=
\frac{\pi}{2}.
}
\]

%%%%%%%%%%%%%%%%%%%%%%%%%%%%%%%%%%%%%%%%%%%%%%%%%%%%%%%%%%%%
\subsection{Critical Delay}
\label{subsec:critical-delay-pure-feedback}
%%%%%%%%%%%%%%%%%%%%%%%%%%%%%%%%%%%%%%%%%%%%%%%%%%%%%%%%%%%%

For

\[
x'(t)
=
-kx(t-\tau),
\]

the first critical delay is

\[
\boxed{
\tau_c
=
\frac{
\pi
}{
2k
}.
}
\]

For

\[
\boxed{
0<k\tau<\frac{\pi}{2},
}
\]

the zero equilibrium is asymptotically stable.

At

\[
\boxed{
k\tau=\frac{\pi}{2},
}
\]

a pair of roots lies on the imaginary axis.

For larger delay, instability appears.

%%%%%%%%%%%%%%%%%%%%%%%%%%%%%%%%%%%%%%%%%%%%%%%%%%%%%%%%%%%%
\subsection{Worked Example: Delay-Induced Instability}
\label{subsec:worked-delay-instability}
%%%%%%%%%%%%%%%%%%%%%%%%%%%%%%%%%%%%%%%%%%%%%%%%%%%%%%%%%%%%

\begin{workedexamplebox}{Critical Delay}

Consider

\[
x'(t)
=
-2x(t-\tau).
\]

Here,

\[
k=2.
\]

The critical delay is

\[
\tau_c
=
\frac{
\pi
}{
2(2)
}
=
\frac{\pi}{4}.
\]

Therefore,

\begin{answerbox}
\[
\boxed{
0<\tau<\frac{\pi}{4}
}
\]

gives asymptotic stability for the zero equilibrium, while the first
stability boundary occurs at

\[
\boxed{
\tau=\frac{\pi}{4}.
}
\]
\end{answerbox}

\end{workedexamplebox}

%%%%%%%%%%%%%%%%%%%%%%%%%%%%%%%%%%%%%%%%%%%%%%%%%%%%%%%%%%%%
\subsection{General Imaginary-Root Conditions}
\label{subsec:general-dde-imaginary-root-conditions}
%%%%%%%%%%%%%%%%%%%%%%%%%%%%%%%%%%%%%%%%%%%%%%%%%%%%%%%%%%%%

For

\[
x'
=
ax+bx(t-\tau),
\]

the characteristic equation is

\[
\lambda-a-be^{-\lambda\tau}=0.
\]

Set

\[
\lambda=i\omega.
\]

Then

\[
i\omega
-
a
-
b
[
\cos(\omega\tau)
-
i\sin(\omega\tau)
]
=
0.
\]

The real part gives

\[
\boxed{
a
+
b\cos(\omega\tau)
=
0,
}
\]

while the imaginary part gives

\[
\boxed{
\omega
+
b\sin(\omega\tau)
=
0,
}
\]

with signs depending on the chosen characteristic-equation arrangement.
The equations must be derived consistently from the original model.

%%%%%%%%%%%%%%%%%%%%%%%%%%%%%%%%%%%%%%%%%%%%%%%%%%%%%%%%%%%%
\subsection{Oscillatory Instability}
\label{subsec:dde-oscillatory-instability}
%%%%%%%%%%%%%%%%%%%%%%%%%%%%%%%%%%%%%%%%%%%%%%%%%%%%%%%%%%%%

When a complex-conjugate pair crosses from the left half-plane into the
right half-plane, the equilibrium loses stability through an oscillatory
mode.

This is fundamentally different from a real eigenvalue crossing zero.

The delay introduces phase lag, which can turn negative feedback into
oscillation.

%%%%%%%%%%%%%%%%%%%%%%%%%%%%%%%%%%%%%%%%%%%%%%%%%%%%%%%%%%%%
\subsection{Delay as Phase Lag}
\label{subsec:delay-phase-lag}
%%%%%%%%%%%%%%%%%%%%%%%%%%%%%%%%%%%%%%%%%%%%%%%%%%%%%%%%%%%%

For a mode

\[
e^{i\omega t},
\]

the delayed state is

\[
e^{i\omega(t-\tau)}
=
e^{-i\omega\tau}
e^{i\omega t}.
\]

Thus the delay contributes a phase shift

\[
\boxed{
-\omega\tau.
}
\]

Large phase lag can make stabilizing feedback act effectively in the wrong
direction.

%%%%%%%%%%%%%%%%%%%%%%%%%%%%%%%%%%%%%%%%%%%%%%%%%%%%%%%%%%%%
\subsection{Hopf Bifurcation in DDEs}
\label{subsec:dde-hopf-bifurcation}
%%%%%%%%%%%%%%%%%%%%%%%%%%%%%%%%%%%%%%%%%%%%%%%%%%%%%%%%%%%%

A delay can act as a bifurcation parameter.

Suppose a complex-conjugate pair of characteristic roots satisfies

\[
\boxed{
\lambda(\tau_c)
=
\pm i\omega_c
}
\]

at a critical delay.

If the roots cross the imaginary axis transversely and appropriate
nondegeneracy conditions hold, a Hopf bifurcation may occur.

A periodic solution can then emerge near the equilibrium.

%%%%%%%%%%%%%%%%%%%%%%%%%%%%%%%%%%%%%%%%%%%%%%%%%%%%%%%%%%%%
\subsection{Delay-Induced Oscillation}
\label{subsec:delay-induced-oscillation}
%%%%%%%%%%%%%%%%%%%%%%%%%%%%%%%%%%%%%%%%%%%%%%%%%%%%%%%%%%%%

One important modeling principle is

\[
\boxed{
\text{stable instantaneous feedback}
+
\text{sufficient delay}
\rightarrow
\text{oscillation or instability}.
}
\]

This occurs in:

\[
\boxed{
\text{control systems},
}
\]

\[
\boxed{
\text{physiological regulation},
}
\]

\[
\boxed{
\text{population dynamics},
}
\]

and

\[
\boxed{
\text{communication networks}.
}
\]

%%%%%%%%%%%%%%%%%%%%%%%%%%%%%%%%%%%%%%%%%%%%%%%%%%%%%%%%%%%%
\subsection{Lambert \(W\) Function}
\label{subsec:dde-lambert-w}
%%%%%%%%%%%%%%%%%%%%%%%%%%%%%%%%%%%%%%%%%%%%%%%%%%%%%%%%%%%%

The Lambert \(W\) function is defined implicitly by

\[
\boxed{
W(z)e^{W(z)}=z.
}
\]

It can express characteristic roots of some simple delay equations.

%%%%%%%%%%%%%%%%%%%%%%%%%%%%%%%%%%%%%%%%%%%%%%%%%%%%%%%%%%%%
\subsection{Lambert \(W\) Solution of a Characteristic Equation}
\label{subsec:lambert-w-characteristic-roots}
%%%%%%%%%%%%%%%%%%%%%%%%%%%%%%%%%%%%%%%%%%%%%%%%%%%%%%%%%%%%

Consider

\[
x'(t)=ax(t)+bx(t-\tau).
\]

The characteristic equation is

\[
\lambda-a
=
be^{-\lambda\tau}.
\]

Multiply by

\[
e^{(\lambda-a)\tau}.
\]

Then

\[
(\lambda-a)
e^{(\lambda-a)\tau}
=
b
e^{-a\tau}.
\]

Multiply by \(\tau\):

\[
(\lambda-a)\tau
e^{(\lambda-a)\tau}
=
b\tau e^{-a\tau}.
\]

Therefore,

\[
\boxed{
(\lambda-a)\tau
=
W_k
\left(
b\tau e^{-a\tau}
\right),
}
\]

where \(W_k\) denotes branch \(k\).

Hence,

\[
\boxed{
\lambda_k
=
a
+
\frac1\tau
W_k
\left(
b\tau e^{-a\tau}
\right).
}
\]

%%%%%%%%%%%%%%%%%%%%%%%%%%%%%%%%%%%%%%%%%%%%%%%%%%%%%%%%%%%%
\subsection{Branches of Lambert \(W\)}
\label{subsec:lambert-w-branches}
%%%%%%%%%%%%%%%%%%%%%%%%%%%%%%%%%%%%%%%%%%%%%%%%%%%%%%%%%%%%

The Lambert \(W\) function has multiple complex branches.

Therefore,

\[
\boxed{
k\in\mathbb{Z}
}
\]

produces infinitely many characteristic roots.

This representation makes the infinite spectrum of the delay equation
explicit.

%%%%%%%%%%%%%%%%%%%%%%%%%%%%%%%%%%%%%%%%%%%%%%%%%%%%%%%%%%%%
\subsection{Nonlinear Delay Equations}
\label{subsec:nonlinear-delay-equations}
%%%%%%%%%%%%%%%%%%%%%%%%%%%%%%%%%%%%%%%%%%%%%%%%%%%%%%%%%%%%

A nonlinear DDE may have form

\[
\boxed{
x'(t)
=
f
\left(
x(t),
x(t-\tau)
\right).
}
\]

As with nonlinear ODEs, exact solutions are uncommon.

Typical analysis focuses on:

\[
\boxed{
\text{equilibria},
}
\]

\[
\boxed{
\text{linear stability},
}
\]

\[
\boxed{
\text{periodic solutions},
}
\]

and

\[
\boxed{
\text{bifurcations}.
}
\]

%%%%%%%%%%%%%%%%%%%%%%%%%%%%%%%%%%%%%%%%%%%%%%%%%%%%%%%%%%%%
\subsection{Linearization of a Nonlinear DDE}
\label{subsec:linearization-nonlinear-dde}
%%%%%%%%%%%%%%%%%%%%%%%%%%%%%%%%%%%%%%%%%%%%%%%%%%%%%%%%%%%%

Suppose

\[
x^*
\]

is an equilibrium of

\[
x'(t)
=
f
\left(
x(t),
x(t-\tau)
\right).
\]

Write

\[
x(t)=x^*+u(t).
\]

Then, to first order,

\[
\boxed{
u'(t)
=
a u(t)
+
b u(t-\tau),
}
\]

where

\[
\boxed{
a
=
\frac{\partial f}{\partial x_1}
(x^*,x^*),
}
\]

and

\[
\boxed{
b
=
\frac{\partial f}{\partial x_2}
(x^*,x^*).
}
\]

The corresponding characteristic equation is

\[
\boxed{
\lambda
=
a
+
be^{-\lambda\tau}.
}
\]

%%%%%%%%%%%%%%%%%%%%%%%%%%%%%%%%%%%%%%%%%%%%%%%%%%%%%%%%%%%%
\subsection{Delayed Logistic Equation}
\label{subsec:delayed-logistic-equation}
%%%%%%%%%%%%%%%%%%%%%%%%%%%%%%%%%%%%%%%%%%%%%%%%%%%%%%%%%%%%

A delayed logistic model is

\[
\boxed{
P'(t)
=
rP(t)
\left[
1-
\frac{
P(t-\tau)
}{
K
}
\right].
}
\]

Here:

\[
r>0
=
\text{growth rate},
\]

\[
K>0
=
\text{carrying capacity},
\]

and

\[
\tau
=
\text{delay in density feedback}.
\]

%%%%%%%%%%%%%%%%%%%%%%%%%%%%%%%%%%%%%%%%%%%%%%%%%%%%%%%%%%%%
\subsection{Equilibria of the Delayed Logistic Model}
\label{subsec:delayed-logistic-equilibria}
%%%%%%%%%%%%%%%%%%%%%%%%%%%%%%%%%%%%%%%%%%%%%%%%%%%%%%%%%%%%

Set

\[
P(t)=P^*.
\]

Then

\[
0
=
rP^*
\left(
1-\frac{P^*}{K}
\right).
\]

Therefore,

\[
\boxed{
P^*=0
}
\]

and

\[
\boxed{
P^*=K.
}
\]

The delay does not change the equilibrium values, but it changes their
stability.

%%%%%%%%%%%%%%%%%%%%%%%%%%%%%%%%%%%%%%%%%%%%%%%%%%%%%%%%%%%%
\subsection{Linearization at Carrying Capacity}
\label{subsec:delayed-logistic-linearization}
%%%%%%%%%%%%%%%%%%%%%%%%%%%%%%%%%%%%%%%%%%%%%%%%%%%%%%%%%%%%

Let

\[
P(t)=K+u(t).
\]

Substitute into

\[
P'
=
rP(t)
\left[
1-\frac{P(t-\tau)}{K}
\right].
\]

Ignoring nonlinear perturbation terms gives

\[
\boxed{
u'(t)
=
-r u(t-\tau).
}
\]

Thus the stability condition is inherited from the pure delayed negative
feedback equation.

%%%%%%%%%%%%%%%%%%%%%%%%%%%%%%%%%%%%%%%%%%%%%%%%%%%%%%%%%%%%
\subsection{Delayed Logistic Stability Threshold}
\label{subsec:delayed-logistic-stability-threshold}
%%%%%%%%%%%%%%%%%%%%%%%%%%%%%%%%%%%%%%%%%%%%%%%%%%%%%%%%%%%%

The linearized equation is

\[
u'
=
-r u(t-\tau).
\]

Therefore the carrying-capacity equilibrium is locally asymptotically
stable when

\[
\boxed{
r\tau
<
\frac{\pi}{2}.
}
\]

The first critical threshold occurs at

\[
\boxed{
r\tau
=
\frac{\pi}{2}.
}
\]

Larger delays can create sustained population oscillations or instability.

%%%%%%%%%%%%%%%%%%%%%%%%%%%%%%%%%%%%%%%%%%%%%%%%%%%%%%%%%%%%
\subsection{Delayed Feedback Control}
\label{subsec:delayed-feedback-control}
%%%%%%%%%%%%%%%%%%%%%%%%%%%%%%%%%%%%%%%%%%%%%%%%%%%%%%%%%%%%

Consider a controlled state

\[
x'(t)
=
ax(t)
+
ku(t),
\]

with delayed feedback

\[
\boxed{
u(t)
=
-x(t-\tau).
}
\]

Then

\[
\boxed{
x'(t)
=
ax(t)
-
k x(t-\tau).
}
\]

Even if the feedback gain \(k\) stabilizes the zero-delay system, a
sufficiently large delay can destabilize it.

%%%%%%%%%%%%%%%%%%%%%%%%%%%%%%%%%%%%%%%%%%%%%%%%%%%%%%%%%%%%
\subsection{Physiological Delay Models}
\label{subsec:physiological-delay-models}
%%%%%%%%%%%%%%%%%%%%%%%%%%%%%%%%%%%%%%%%%%%%%%%%%%%%%%%%%%%%

Physiological regulation often involves delays caused by:

\[
\boxed{
\text{hormone production},
}
\]

\[
\boxed{
\text{blood transport},
}
\]

\[
\boxed{
\text{cell maturation},
}
\]

or

\[
\boxed{
\text{neural response}.
}
\]

Feedback models with delays can exhibit oscillations even when their
instantaneous counterparts converge monotonically.

%%%%%%%%%%%%%%%%%%%%%%%%%%%%%%%%%%%%%%%%%%%%%%%%%%%%%%%%%%%%
\subsection{Epidemiological Delay Models}
\label{subsec:epidemiological-delay-models}
%%%%%%%%%%%%%%%%%%%%%%%%%%%%%%%%%%%%%%%%%%%%%%%%%%%%%%%%%%%%

Delays can represent:

\[
\boxed{
\text{incubation time},
}
\]

\[
\boxed{
\text{latent periods},
}
\]

\[
\boxed{
\text{treatment delays},
}
\]

or

\[
\boxed{
\text{temporary immunity}.
}
\]

For example, infection terms can depend on population states at earlier
times rather than only at the present instant.

%%%%%%%%%%%%%%%%%%%%%%%%%%%%%%%%%%%%%%%%%%%%%%%%%%%%%%%%%%%%
\subsection{Maturation Delay}
\label{subsec:maturation-delay}
%%%%%%%%%%%%%%%%%%%%%%%%%%%%%%%%%%%%%%%%%%%%%%%%%%%%%%%%%%%%

Suppose individuals require time \(\tau\) to mature.

Recruitment into the adult population at time \(t\) may depend on births
at time

\[
t-\tau.
\]

A schematic model is

\[
\boxed{
A'(t)
=
bA(t-\tau)
-
dA(t).
}
\]

The delayed term represents recruitment from the earlier cohort.

%%%%%%%%%%%%%%%%%%%%%%%%%%%%%%%%%%%%%%%%%%%%%%%%%%%%%%%%%%%%
\subsection{Systems of Delay Differential Equations}
\label{subsec:systems-of-ddes}
%%%%%%%%%%%%%%%%%%%%%%%%%%%%%%%%%%%%%%%%%%%%%%%%%%%%%%%%%%%%

A vector DDE may be written

\[
\boxed{
\mathbf{x}'(t)
=
\mathbf{f}
\left(
\mathbf{x}(t),
\mathbf{x}(t-\tau)
\right).
}
\]

A linear system has form

\[
\boxed{
\mathbf{x}'(t)
=
A\mathbf{x}(t)
+
B\mathbf{x}(t-\tau).
}
\]

%%%%%%%%%%%%%%%%%%%%%%%%%%%%%%%%%%%%%%%%%%%%%%%%%%%%%%%%%%%%
\subsection{Characteristic Equation for Linear DDE Systems}
\label{subsec:linear-dde-system-characteristic}
%%%%%%%%%%%%%%%%%%%%%%%%%%%%%%%%%%%%%%%%%%%%%%%%%%%%%%%%%%%%

Try

\[
\mathbf{x}(t)
=
\mathbf{v}e^{\lambda t}.
\]

Then

\[
\lambda\mathbf{v}
=
A\mathbf{v}
+
Be^{-\lambda\tau}\mathbf{v}.
\]

Thus nontrivial \(\mathbf{v}\) requires

\[
\boxed{
\det
\left[
\lambda I
-
A
-
Be^{-\lambda\tau}
\right]
=
0.
}
\]

This is the matrix-valued delay characteristic equation.

%%%%%%%%%%%%%%%%%%%%%%%%%%%%%%%%%%%%%%%%%%%%%%%%%%%%%%%%%%%%
\subsection{Several Delays in a Linear System}
\label{subsec:linear-system-multiple-delays}
%%%%%%%%%%%%%%%%%%%%%%%%%%%%%%%%%%%%%%%%%%%%%%%%%%%%%%%%%%%%

For

\[
\mathbf{x}'(t)
=
A_0\mathbf{x}(t)
+
\sum_{j=1}^{m}
A_j
\mathbf{x}(t-\tau_j),
\]

the characteristic equation is

\[
\boxed{
\det
\left[
\lambda I
-
A_0
-
\sum_{j=1}^{m}
A_je^{-\lambda\tau_j}
\right]
=
0.
}
\]

The spectral problem can become highly complicated.

%%%%%%%%%%%%%%%%%%%%%%%%%%%%%%%%%%%%%%%%%%%%%%%%%%%%%%%%%%%%
\subsection{Neutral Delay Differential Equations}
\label{subsec:neutral-delay-differential-equations}
%%%%%%%%%%%%%%%%%%%%%%%%%%%%%%%%%%%%%%%%%%%%%%%%%%%%%%%%%%%%

A neutral DDE includes delayed derivatives.

For example,

\[
\boxed{
x'(t)
-
c x'(t-\tau)
=
ax(t)
+
bx(t-\tau).
}
\]

This differs fundamentally from a retarded DDE.

%%%%%%%%%%%%%%%%%%%%%%%%%%%%%%%%%%%%%%%%%%%%%%%%%%%%%%%%%%%%
\subsection{Retarded Versus Neutral DDEs}
\label{subsec:retarded-versus-neutral-ddes}
%%%%%%%%%%%%%%%%%%%%%%%%%%%%%%%%%%%%%%%%%%%%%%%%%%%%%%%%%%%%

A retarded DDE involves delayed states:

\[
\boxed{
x(t-\tau).
}
\]

A neutral DDE may also involve delayed derivatives:

\[
\boxed{
x'(t-\tau).
}
\]

Neutral equations can have more delicate stability and regularity
properties.

%%%%%%%%%%%%%%%%%%%%%%%%%%%%%%%%%%%%%%%%%%%%%%%%%%%%%%%%%%%%
\subsection{Characteristic Equation for a Neutral DDE}
\label{subsec:neutral-dde-characteristic}
%%%%%%%%%%%%%%%%%%%%%%%%%%%%%%%%%%%%%%%%%%%%%%%%%%%%%%%%%%%%

For

\[
x'(t)
-
c x'(t-\tau)
=
ax(t)
+
bx(t-\tau),
\]

try

\[
x=e^{\lambda t}.
\]

Then

\[
\lambda
-
c\lambda e^{-\lambda\tau}
=
a
+
be^{-\lambda\tau}.
\]

Thus,

\[
\boxed{
\lambda
\left(
1-ce^{-\lambda\tau}
\right)
-
a
-
be^{-\lambda\tau}
=
0.
}
\]

%%%%%%%%%%%%%%%%%%%%%%%%%%%%%%%%%%%%%%%%%%%%%%%%%%%%%%%%%%%%
\subsection{Distributed Delay as an Integral Operator}
\label{subsec:distributed-delay-integral-operator}
%%%%%%%%%%%%%%%%%%%%%%%%%%%%%%%%%%%%%%%%%%%%%%%%%%%%%%%%%%%%

Consider

\[
x'(t)
=
ax(t)
+
\int_0^\tau
K(s)x(t-s)\,ds.
\]

The delayed history acts through an integral operator.

Thus distributed-delay systems sit naturally between:

\[
\boxed{
\text{DDEs}
}
\]

and

\[
\boxed{
\text{Volterra integral equations}.
}
\]

%%%%%%%%%%%%%%%%%%%%%%%%%%%%%%%%%%%%%%%%%%%%%%%%%%%%%%%%%%%%
\subsection{Characteristic Equation for Distributed Delay}
\label{subsec:distributed-delay-characteristic}
%%%%%%%%%%%%%%%%%%%%%%%%%%%%%%%%%%%%%%%%%%%%%%%%%%%%%%%%%%%%

Try

\[
x(t)=e^{\lambda t}.
\]

Then

\[
x(t-s)
=
e^{-\lambda s}
e^{\lambda t}.
\]

Therefore,

\[
\lambda
=
a
+
\int_0^\tau
K(s)e^{-\lambda s}\,ds.
\]

Thus,

\[
\boxed{
\lambda
-
a
-
\int_0^\tau
K(s)e^{-\lambda s}\,ds
=
0.
}
\]

The kernel enters through a Laplace-type transform.

%%%%%%%%%%%%%%%%%%%%%%%%%%%%%%%%%%%%%%%%%%%%%%%%%%%%%%%%%%%%
\subsection{Exponential Memory Kernel}
\label{subsec:exponential-memory-kernel}
%%%%%%%%%%%%%%%%%%%%%%%%%%%%%%%%%%%%%%%%%%%%%%%%%%%%%%%%%%%%

Suppose

\[
K(s)
=
\alpha e^{-\beta s}.
\]

Then the memory contribution decays exponentially with lag.

Such kernels can sometimes be converted into additional ODE state
variables, replacing an infinite-memory integral by a finite-dimensional
augmented system.

This technique is often called the linear-chain or auxiliary-variable
idea.

%%%%%%%%%%%%%%%%%%%%%%%%%%%%%%%%%%%%%%%%%%%%%%%%%%%%%%%%%%%%
\subsection{State Augmentation for Exponential Memory}
\label{subsec:state-augmentation-memory}
%%%%%%%%%%%%%%%%%%%%%%%%%%%%%%%%%%%%%%%%%%%%%%%%%%%%%%%%%%%%

Define

\[
\boxed{
z(t)
=
\int_0^\infty
\alpha e^{-\beta s}
x(t-s)\,ds.
}
\]

Under suitable assumptions, differentiation yields an equation of the
form

\[
\boxed{
z'
=
\alpha x
-
\beta z.
}
\]

Thus a distributed-memory problem can sometimes be rewritten as a larger
ODE system.

%%%%%%%%%%%%%%%%%%%%%%%%%%%%%%%%%%%%%%%%%%%%%%%%%%%%%%%%%%%%
\subsection{DDEs and Integral Equations}
\label{subsec:dde-integral-equation-connection}
%%%%%%%%%%%%%%%%%%%%%%%%%%%%%%%%%%%%%%%%%%%%%%%%%%%%%%%%%%%%

Integrating

\[
x'(t)
=
f
\left(
t,
x(t),
x(t-\tau)
\right)
\]

from \(0\) to \(t\) gives

\[
\boxed{
x(t)
=
x(0)
+
\int_0^t
f
\left(
s,
x(s),
x(s-\tau)
\right)
\,ds.
}
\]

Thus a DDE can also be viewed as an integral equation with delayed
arguments.

%%%%%%%%%%%%%%%%%%%%%%%%%%%%%%%%%%%%%%%%%%%%%%%%%%%%%%%%%%%%
\subsection{Laplace Transform of a Delayed Function}
\label{subsec:laplace-delayed-function-dde}
%%%%%%%%%%%%%%%%%%%%%%%%%%%%%%%%%%%%%%%%%%%%%%%%%%%%%%%%%%%%

When a function is zero before \(t=0\), the standard delay property is

\[
\boxed{
\mathcal{L}
\{
u_\tau(t)f(t-\tau)
\}
=
e^{-\tau s}
F(s).
}
\]

However, DDEs usually have nonzero history for negative times.

Therefore the Laplace transform of

\[
x(t-\tau)
\]

must include contributions from the prescribed history.

%%%%%%%%%%%%%%%%%%%%%%%%%%%%%%%%%%%%%%%%%%%%%%%%%%%%%%%%%%%%
\subsection{Laplace Transform with a History Function}
\label{subsec:laplace-history-function}
%%%%%%%%%%%%%%%%%%%%%%%%%%%%%%%%%%%%%%%%%%%%%%%%%%%%%%%%%%%%

For \(t\geq0\),

\[
\mathcal{L}\{x(t-\tau)\}
=
\int_0^\infty
e^{-st}x(t-\tau)\,dt.
\]

Let

\[
u=t-\tau.
\]

Then

\[
\mathcal{L}\{x(t-\tau)\}
=
e^{-s\tau}
\int_{-\tau}^{\infty}
e^{-su}x(u)\,du.
\]

Split the integral:

\[
\boxed{
\mathcal{L}\{x(t-\tau)\}
=
e^{-s\tau}
\left[
\int_{-\tau}^{0}
e^{-su}\phi(u)\,du
+
X(s)
\right].
}
\]

Thus history terms appear explicitly.

%%%%%%%%%%%%%%%%%%%%%%%%%%%%%%%%%%%%%%%%%%%%%%%%%%%%%%%%%%%%
\subsection{Worked Example: Laplace Transform with Constant History}
\label{subsec:worked-laplace-dde-history}
%%%%%%%%%%%%%%%%%%%%%%%%%%%%%%%%%%%%%%%%%%%%%%%%%%%%%%%%%%%%

\begin{workedexamplebox}{Transforming a Delayed State}

Suppose

\[
x(t)=1,
\qquad
-1\leq t\leq0.
\]

Then

\[
\mathcal{L}\{x(t-1)\}
=
e^{-s}
\left[
X(s)
+
\int_{-1}^{0}
e^{-su}\,du
\right].
\]

The history integral is

\[
\int_{-1}^{0}
e^{-su}\,du
=
\frac{
e^s-1
}{
s
}.
\]

Therefore,

\begin{answerbox}
\[
\boxed{
\mathcal{L}\{x(t-1)\}
=
e^{-s}X(s)
+
\frac{
1-e^{-s}
}{
s
}.
}
\]
\end{answerbox}

\end{workedexamplebox}

%%%%%%%%%%%%%%%%%%%%%%%%%%%%%%%%%%%%%%%%%%%%%%%%%%%%%%%%%%%%
\subsection{Numerical Solution of DDEs}
\label{subsec:numerical-solution-ddes}
%%%%%%%%%%%%%%%%%%%%%%%%%%%%%%%%%%%%%%%%%%%%%%%%%%%%%%%%%%%%

Most nonlinear DDEs cannot be solved analytically.

Numerical DDE solvers extend ODE methods while maintaining access to past
solution values.

At time \(t_n\), the method may require

\[
\boxed{
x(t_n-\tau),
}
\]

which generally does not coincide with a stored grid point.

%%%%%%%%%%%%%%%%%%%%%%%%%%%%%%%%%%%%%%%%%%%%%%%%%%%%%%%%%%%%
\subsection{Interpolation of Delayed States}
\label{subsec:dde-interpolation}
%%%%%%%%%%%%%%%%%%%%%%%%%%%%%%%%%%%%%%%%%%%%%%%%%%%%%%%%%%%%

If

\[
t_n-\tau
\]

lies between computed solution points, the solver must interpolate the
past solution.

Therefore DDE solvers often maintain a continuous approximation of the
recent solution history.

This is sometimes called a continuous extension or dense output.

%%%%%%%%%%%%%%%%%%%%%%%%%%%%%%%%%%%%%%%%%%%%%%%%%%%%%%%%%%%%
\subsection{Euler Method for a DDE}
\label{subsec:euler-method-dde}
%%%%%%%%%%%%%%%%%%%%%%%%%%%%%%%%%%%%%%%%%%%%%%%%%%%%%%%%%%%%

Consider

\[
x'(t)
=
f
\left(
t,
x(t),
x(t-\tau)
\right).
\]

With time step \(h\),

\[
\boxed{
x_{n+1}
=
x_n
+
h
f
\left(
t_n,
x_n,
x(t_n-\tau)
\right).
}
\]

The delayed value must come from:

\[
\boxed{
\text{the history function},
}
\]

or

\[
\boxed{
\text{previous numerical solution values}.
}
\]

%%%%%%%%%%%%%%%%%%%%%%%%%%%%%%%%%%%%%%%%%%%%%%%%%%%%%%%%%%%%
\subsection{Delay Aligned with the Numerical Grid}
\label{subsec:dde-grid-aligned-delay}
%%%%%%%%%%%%%%%%%%%%%%%%%%%%%%%%%%%%%%%%%%%%%%%%%%%%%%%%%%%%

If

\[
\tau=mh
\]

for an integer \(m\), then

\[
t_n-\tau
=
t_{n-m}.
\]

The Euler update becomes

\[
\boxed{
x_{n+1}
=
x_n
+
h
f
\left(
t_n,
x_n,
x_{n-m}
\right).
}
\]

This avoids interpolation.

%%%%%%%%%%%%%%%%%%%%%%%%%%%%%%%%%%%%%%%%%%%%%%%%%%%%%%%%%%%%
\subsection{Worked Example: Euler DDE Step}
\label{subsec:worked-euler-dde}
%%%%%%%%%%%%%%%%%%%%%%%%%%%%%%%%%%%%%%%%%%%%%%%%%%%%%%%%%%%%

\begin{workedexamplebox}{Delayed Euler Approximation}

Consider

\[
x'(t)
=
-x(t-1),
\]

with history

\[
x(t)=1,
\qquad
-1\leq t\leq0.
\]

Use

\[
h=0.25.
\]

At \(t_0=0\),

\[
x_0=1,
\]

and

\[
x(-1)=1.
\]

Thus,

\[
x_1
=
x_0
+
0.25[-x(-1)].
\]

Therefore,

\begin{answerbox}
\[
\boxed{
x_1
=
1-0.25
=
0.75.
}
\]
\end{answerbox}

\end{workedexamplebox}

%%%%%%%%%%%%%%%%%%%%%%%%%%%%%%%%%%%%%%%%%%%%%%%%%%%%%%%%%%%%
\subsection{Discontinuity Tracking}
\label{subsec:dde-discontinuity-tracking}
%%%%%%%%%%%%%%%%%%%%%%%%%%%%%%%%%%%%%%%%%%%%%%%%%%%%%%%%%%%%

If the history and DDE are not fully compatible at \(t=0\), derivative
discontinuities can occur.

These discontinuities propagate to times such as

\[
\boxed{
\tau,
\ 2\tau,
\ 3\tau,\ldots
}
\]

and may increase in derivative order as the solution evolves.

Accurate DDE solvers often track these points explicitly.

%%%%%%%%%%%%%%%%%%%%%%%%%%%%%%%%%%%%%%%%%%%%%%%%%%%%%%%%%%%%
\subsection{Numerical Stability}
\label{subsec:dde-numerical-stability}
%%%%%%%%%%%%%%%%%%%%%%%%%%%%%%%%%%%%%%%%%%%%%%%%%%%%%%%%%%%%

Numerical stability is especially important because the true DDE can be
near a delay-induced stability boundary.

A method may:

\[
\boxed{
\text{create artificial oscillation},
}
\]

\[
\boxed{
\text{over-damp true oscillation},
}
\]

or

\[
\boxed{
\text{shift a critical delay}.
}
\]

Step-size convergence should therefore be checked.

%%%%%%%%%%%%%%%%%%%%%%%%%%%%%%%%%%%%%%%%%%%%%%%%%%%%%%%%%%%%
\subsection{Bifurcation Diagrams for DDEs}
\label{subsec:dde-bifurcation-diagrams}
%%%%%%%%%%%%%%%%%%%%%%%%%%%%%%%%%%%%%%%%%%%%%%%%%%%%%%%%%%%%

A parameter such as

\[
\boxed{
\tau
}
\]

can be varied while tracking:

\[
\boxed{
\text{equilibria},
}
\]

\[
\boxed{
\text{characteristic roots},
}
\]

\[
\boxed{
\text{periodic solutions},
}
\]

and

\[
\boxed{
\text{oscillation amplitudes}.
}
\]

This reveals delay-induced bifurcations.

%%%%%%%%%%%%%%%%%%%%%%%%%%%%%%%%%%%%%%%%%%%%%%%%%%%%%%%%%%%%
\subsection{DDE Modeling Workflow}
\label{subsec:dde-modeling-workflow}
%%%%%%%%%%%%%%%%%%%%%%%%%%%%%%%%%%%%%%%%%%%%%%%%%%%%%%%%%%%%

A practical workflow is:

\begin{enumerate}
    \item identify the process creating the delay;
    \item determine whether the delay is fixed, distributed, or
          state-dependent;
    \item determine the maximum delay;
    \item specify the required history function;
    \item identify equilibria;
    \item linearize around equilibria;
    \item derive the characteristic equation;
    \item analyze characteristic roots;
    \item locate imaginary-axis crossings;
    \item identify possible Hopf bifurcations;
    \item solve by the method of steps when possible;
    \item use numerical DDE solvers for nonlinear models;
    \item vary the delay and other parameters;
    \item compare delayed and zero-delay behavior;
    \item interpret oscillations and instability in the application.
\end{enumerate}

%%%%%%%%%%%%%%%%%%%%%%%%%%%%%%%%%%%%%%%%%%%%%%%%%%%%%%%%%%%%
\subsection{Choosing a DDE Method}
\label{subsec:choosing-dde-method}
%%%%%%%%%%%%%%%%%%%%%%%%%%%%%%%%%%%%%%%%%%%%%%%%%%%%%%%%%%%%

For simple constant-delay equations with known history, consider:

\[
\boxed{
\text{method of steps}.
}
\]

For linear stability, use:

\[
\boxed{
\text{characteristic roots}.
}
\]

For simple scalar characteristic equations, consider:

\[
\boxed{
\text{Lambert }W.
}
\]

For convolution-type memory, consider:

\[
\boxed{
\text{Laplace transforms}.
}
\]

For distributed exponential memory, consider:

\[
\boxed{
\text{state augmentation}.
}
\]

For nonlinear models, use:

\[
\boxed{
\text{linearization and numerical continuation}.
}
\]

%%%%%%%%%%%%%%%%%%%%%%%%%%%%%%%%%%%%%%%%%%%%%%%%%%%%%%%%%%%%
\subsection{Common Errors}
\label{subsec:dde-common-errors}
%%%%%%%%%%%%%%%%%%%%%%%%%%%%%%%%%%%%%%%%%%%%%%%%%%%%%%%%%%%%

\subsubsection{Providing Only \(x(0)\)}

A DDE generally requires a history function over an interval such as

\[
[-\tau,0].
\]

One scalar initial value is usually insufficient.

%%%%%%%%%%%%%%%%%%%%%%%%%%%%%%%%%%%%%%%%%%%%%%%%%%%%%%%%%%%%
\subsubsection{Treating \(x(t-\tau)\) as \(x(t)-\tau\)}

These expressions are completely different:

\[
\boxed{
x(t-\tau)
}
\]

is the function evaluated at an earlier time.

It is not

\[
\boxed{
x(t)-\tau.
}
\]

%%%%%%%%%%%%%%%%%%%%%%%%%%%%%%%%%%%%%%%%%%%%%%%%%%%%%%%%%%%%
\subsubsection{Assuming Delay Changes the Equilibrium Location}

For many autonomous equations, a constant equilibrium satisfies

\[
x(t-\tau)=x(t)=x^*.
\]

Thus delay often changes stability rather than equilibrium location.

%%%%%%%%%%%%%%%%%%%%%%%%%%%%%%%%%%%%%%%%%%%%%%%%%%%%%%%%%%%%
\subsubsection{Using a Polynomial Characteristic Equation}

The delay characteristic equation contains

\[
e^{-\lambda\tau}.
\]

It is transcendental and generally has infinitely many roots.

%%%%%%%%%%%%%%%%%%%%%%%%%%%%%%%%%%%%%%%%%%%%%%%%%%%%%%%%%%%%
\subsubsection{Checking Only One Characteristic Root}

Stability requires control of the entire spectrum.

A single stable root does not imply the DDE is stable.

%%%%%%%%%%%%%%%%%%%%%%%%%%%%%%%%%%%%%%%%%%%%%%%%%%%%%%%%%%%%
\subsubsection{Ignoring Complex Roots}

Delay-induced instability often occurs through a complex-conjugate pair.

Restricting attention to real roots can miss the critical behavior.

%%%%%%%%%%%%%%%%%%%%%%%%%%%%%%%%%%%%%%%%%%%%%%%%%%%%%%%%%%%%
\subsubsection{Using Zero-Delay Stability as the Final Answer}

A system that is stable for

\[
\tau=0
\]

may become unstable for positive delay.

The delay must be analyzed explicitly.

%%%%%%%%%%%%%%%%%%%%%%%%%%%%%%%%%%%%%%%%%%%%%%%%%%%%%%%%%%%%
\subsubsection{Using the Ordinary Laplace Delay Rule Without History Terms}

For a DDE with nonzero negative-time history,

\[
\mathcal{L}\{x(t-\tau)\}
\]

contains an extra integral involving the history function.

%%%%%%%%%%%%%%%%%%%%%%%%%%%%%%%%%%%%%%%%%%%%%%%%%%%%%%%%%%%%
\subsubsection{Assuming Every Delay Produces Instability}

Small delays may preserve stability.

The effect depends on the coefficients and feedback structure.

%%%%%%%%%%%%%%%%%%%%%%%%%%%%%%%%%%%%%%%%%%%%%%%%%%%%%%%%%%%%
\subsubsection{Confusing Retarded and Neutral Equations}

A retarded DDE uses delayed states.

A neutral DDE contains delayed derivatives and requires separate theory.

%%%%%%%%%%%%%%%%%%%%%%%%%%%%%%%%%%%%%%%%%%%%%%%%%%%%%%%%%%%%
\subsubsection{Ignoring State-Dependent Delay Constraints}

For

\[
t-\tau(x(t)),
\]

the delayed argument must remain in the region where history or solution
data are available.

%%%%%%%%%%%%%%%%%%%%%%%%%%%%%%%%%%%%%%%%%%%%%%%%%%%%%%%%%%%%
\subsubsection{Ignoring Interpolation Error in Numerical Solvers}

If the delayed time does not lie on the numerical mesh, interpolation
becomes part of the numerical approximation.

Its accuracy matters.

%%%%%%%%%%%%%%%%%%%%%%%%%%%%%%%%%%%%%%%%%%%%%%%%%%%%%%%%%%%%
\subsubsection{Ignoring Propagating Derivative Discontinuities}

A DDE solution can be continuous while derivatives have jumps at
predictable delay-related times.

Numerical methods should respect those discontinuities.

%%%%%%%%%%%%%%%%%%%%%%%%%%%%%%%%%%%%%%%%%%%%%%%%%%%%%%%%%%%%
\subsection{Applications}
\label{subsec:dde-applications}
%%%%%%%%%%%%%%%%%%%%%%%%%%%%%%%%%%%%%%%%%%%%%%%%%%%%%%%%%%%%

\begin{applicationbox}

\textbf{Population Biology.}
Maturation, gestation, recruitment, and delayed density dependence lead
naturally to DDE models.

\medskip

\textbf{Epidemiology.}
Incubation, immunity, treatment, and reporting delays can be represented
through delayed states.

\medskip

\textbf{Physiology.}
Hormonal regulation, blood-cell production, neural feedback, and
respiratory control involve biologically meaningful response delays.

\medskip

\textbf{Control Theory.}
Sensors, processors, actuators, and communication networks introduce
feedback delays that can destabilize otherwise stable systems.

\medskip

\textbf{Engineering.}
Transport systems, machining processes, power networks, and
communication systems often contain unavoidable delays.

\medskip

\textbf{Economics.}
Investment, production, policy, and information delays can create
oscillation and instability.

\medskip

\textbf{Ecology.}
Delayed resource limitation and age structure can produce cycles in
population models.

\medskip

\textbf{Network Science.}
Propagation delays influence synchronization, consensus, and stability in
interconnected systems.

\end{applicationbox}

%%%%%%%%%%%%%%%%%%%%%%%%%%%%%%%%%%%%%%%%%%%%%%%%%%%%%%%%%%%%
\subsection{Exercises}
\label{subsec:dde-exercises}
%%%%%%%%%%%%%%%%%%%%%%%%%%%%%%%%%%%%%%%%%%%%%%%%%%%%%%%%%%%%

\begin{exercise}[1]
Define a delay differential equation.
\end{exercise}

\begin{exercise}[1]
Define a discrete delay.
\end{exercise}

\begin{exercise}[1]
Define a distributed delay.
\end{exercise}

\begin{exercise}[1]
Define a history function.
\end{exercise}

\begin{exercise}[1]
Explain why a scalar DDE is effectively infinite-dimensional.
\end{exercise}

\begin{exercise}[1]
Define a neutral DDE.
\end{exercise}

\begin{exercise}[2]
For

\[
x'(t)
=
x(t-2),
\]

state the history interval needed to begin the solution at \(t=0\).
\end{exercise}

\begin{exercise}[2]
Find the equilibria of

\[
x'(t)
=
x(t)
[
1-x(t-\tau)
].
\]
\end{exercise}

\begin{exercise}[2]
Derive the characteristic equation for

\[
x'(t)
=
3x(t)
-
2x(t-\tau).
\]
\end{exercise}

\begin{exercise}[2]
Derive the characteristic equation for

\[
x'(t)
=
-kx(t-\tau).
\]
\end{exercise}

\begin{exercise}[2]
Find the first critical delay for

\[
x'(t)
=
-4x(t-\tau).
\]
\end{exercise}

\begin{exercise}[2]
Determine whether the equilibrium is stable when

\[
k\tau=\frac{\pi}{4}
\]

for

\[
x'=-kx(t-\tau).
\]
\end{exercise}

\begin{exercise}[3]
Use the method of steps to solve

\[
x'(t)=x(t-1)
\]

on

\[
0\leq t\leq2
\]

when

\[
x(t)=2,
\qquad
-1\leq t\leq0.
\]
\end{exercise}

\begin{exercise}[3]
Solve

\[
x'(t)
=
-x(t-1)
\]

on the first step interval using a constant history

\[
x(t)=1,
\qquad
-1\leq t\leq0.
\]
\end{exercise}

\begin{exercise}[3]
Explain how derivative discontinuities propagate through the method of
steps.
\end{exercise}

\begin{exercise}[3]
Derive the characteristic equation

\[
\lambda
=
a
+
be^{-\lambda\tau}
\]

for the scalar linear DDE.
\end{exercise}

\begin{exercise}[3]
Show that the zero-delay limit of the previous equation is

\[
x'=(a+b)x.
\]
\end{exercise}

\begin{exercise}[3]
For

\[
x'(t)
=
-kx(t-\tau),
\]

derive the imaginary-axis equations using

\[
\lambda=i\omega.
\]
\end{exercise}

\begin{exercise}[3]
Derive the stability threshold

\[
k\tau=\frac{\pi}{2}.
\]
\end{exercise}

\begin{exercise}[3]
Explain physically why delay can destabilize negative feedback.
\end{exercise}

\begin{exercise}[3]
Derive the Lambert \(W\) expression for the roots of

\[
x'=ax+bx(t-\tau).
\]
\end{exercise}

\begin{exercise}[3]
Linearize the delayed logistic equation around

\[
P=K.
\]
\end{exercise}

\begin{exercise}[3]
Derive the delayed logistic stability threshold

\[
r\tau<\frac{\pi}{2}.
\]
\end{exercise}

\begin{exercise}[3]
Formulate a maturation-delay population model.
\end{exercise}

\begin{exercise}[3]
Formulate a delayed feedback-control model.
\end{exercise}

\begin{exercise}[3]
Derive the characteristic equation for

\[
\mathbf{x}'
=
A\mathbf{x}(t)
+
B\mathbf{x}(t-\tau).
\]
\end{exercise}

\begin{exercise}[3]
Derive the characteristic equation for a linear equation with two
distinct delays.
\end{exercise}

\begin{exercise}[3]
Derive the characteristic equation of a neutral DDE.
\end{exercise}

\begin{exercise}[3]
Show how a distributed-delay exponential mode leads to

\[
\int_0^\tau
K(s)e^{-\lambda s}\,ds.
\]
\end{exercise}

\begin{exercise}[3]
Convert a DDE into an integral equation by integrating from \(0\) to
\(t\).
\end{exercise}

\begin{exercise}[3]
Derive the Laplace-transform formula for \(x(t-\tau)\) when a nonzero
history function is prescribed.
\end{exercise}

\begin{exercise}[3]
Construct an Euler approximation for

\[
x'(t)
=
-x(t-1)
\]

using

\[
h=0.25.
\]
\end{exercise}

\begin{exercise}[3]
Explain why interpolation is required when

\[
\tau/h
\]

is not an integer.
\end{exercise}

\begin{exercise}[4]
Study existence and uniqueness theory for retarded functional
differential equations.
\end{exercise}

\begin{exercise}[4]
Study the solution semigroup on

\[
C([-\tau,0],\mathbb{R}^n).
\]
\end{exercise}

\begin{exercise}[4]
Investigate spectral theory of linear delay equations.
\end{exercise}

\begin{exercise}[4]
Prove the stability threshold for pure delayed negative feedback using
the argument principle or root-continuation methods.
\end{exercise}

\begin{exercise}[4]
Study root crossing direction at a delay-induced Hopf bifurcation.
\end{exercise}

\begin{exercise}[4]
Derive Hopf bifurcation conditions for the delayed logistic equation.
\end{exercise}

\begin{exercise}[4]
Study normal-form theory for DDE Hopf bifurcations.
\end{exercise}

\begin{exercise}[4]
Investigate multiple stability switches in equations with several delays.
\end{exercise}

\begin{exercise}[4]
Study the Lambert \(W\) branches associated with DDE spectra.
\end{exercise}

\begin{exercise}[4]
Investigate distributed delays and their Laplace transforms.
\end{exercise}

\begin{exercise}[4]
Study the linear-chain trick for gamma-distributed delays.
\end{exercise}

\begin{exercise}[4]
Study neutral functional differential equations.
\end{exercise}

\begin{exercise}[4]
Investigate state-dependent delays.
\end{exercise}

\begin{exercise}[4]
Study delayed epidemic models.
\end{exercise}

\begin{exercise}[4]
Study delayed predator--prey systems.
\end{exercise}

\begin{exercise}[4]
Analyze delayed synchronization in networks.
\end{exercise}

\begin{exercise}[4]
Study numerical Runge--Kutta methods for DDEs.
\end{exercise}

\begin{exercise}[4]
Investigate interpolation and discontinuity tracking in numerical DDE
solvers.
\end{exercise}

\begin{exercise}[4]
Study numerical bifurcation continuation for periodic DDE solutions.
\end{exercise}

%%%%%%%%%%%%%%%%%%%%%%%%%%%%%%%%%%%%%%%%%%%%%%%%%%%%%%%%%%%%
\subsection{Conceptual Summary}
\label{subsec:dde-summary}
%%%%%%%%%%%%%%%%%%%%%%%%%%%%%%%%%%%%%%%%%%%%%%%%%%%%%%%%%%%%

A delay differential equation has the general form

\[
\boxed{
x'(t)
=
f
\left(
t,
x(t),
x(t-\tau)
\right).
}
\]

Unlike an ODE, a DDE generally requires a history function

\[
\boxed{
x(t)
=
\phi(t),
\qquad
-\tau\leq t\leq0.
}
\]

For a linear scalar DDE,

\[
\boxed{
x'(t)
=
ax(t)
+
bx(t-\tau),
}
\]

the exponential trial

\[
x(t)=e^{\lambda t}
\]

produces the transcendental characteristic equation

\[
\boxed{
\lambda
=
a
+
be^{-\lambda\tau}.
}
\]

Because of the exponential delay factor, there are generally infinitely
many characteristic roots.

For pure delayed negative feedback,

\[
\boxed{
x'(t)
=
-kx(t-\tau),
}
\]

the first stability boundary occurs at

\[
\boxed{
k\tau
=
\frac{\pi}{2}.
}
\]

Thus increasing delay can destabilize an otherwise stable feedback
system.

The state of a DDE is a history segment

\[
\boxed{
x_t(\theta)=x(t+\theta),
}
\]

so DDEs are naturally infinite-dimensional dynamical systems.

Delay differential equations therefore combine

\[
\boxed{
\text{ODEs}
+
\text{memory}
+
\text{functional analysis}
+
\text{stability}
+
\text{bifurcation theory}.
}
\]

%%%%%%%%%%%%%%%%%%%%%%%%%%%%%%%%%%%%%%%%%%%%%%%%%%%%%%%%%%%%
\subsection{Section Connections}
\label{subsec:dde-connections}
%%%%%%%%%%%%%%%%%%%%%%%%%%%%%%%%%%%%%%%%%%%%%%%%%%%%%%%%%%%%

\chapterconnection{
Delay Differential Equations extend ordinary differential equations by
allowing present rates of change to depend on past states. Their history
functions connect naturally with the Volterra and memory formulations of
Section~\ref{sec:integral-equations}, while their characteristic equations
extend the stability analysis of Section~\ref{sec:dynamical-systems}.
Because delayed feedback can generate oscillations and Hopf bifurcations,
DDEs provide an important bridge between modeling, control theory, and
infinite-dimensional dynamical systems. The next section,
Differential-Algebraic Equations, completes Chapter 9 by studying systems
in which differential evolution is constrained simultaneously by
algebraic equations.
}

%%%%%%%%%%%%%%%%%%%%%%%%%%%%%%%%%%%%%%%%%%%%%%%%%%%%%%%%%%%%
% SECTION SOURCE TRACKING
%%%%%%%%%%%%%%%%%%%%%%%%%%%%%%%%%%%%%%%%%%%%%%%%%%%%%%%%%%%%

%------------------------------------------------------------
% 9.7 Delay Differential Equations
%------------------------------------------------------------
%
% Source basis:
%   - Standard delay differential equation theory.
%   - Standard functional differential equation stability theory.
%   - Standard delayed-feedback and mathematical-modeling methods.
%
% Used for:
%   - DDE definitions;
%   - discrete, multiple, time-dependent, and state-dependent delays;
%   - distributed delays;
%   - history functions;
%   - phase space of DDEs;
%   - infinite-dimensional interpretation;
%   - method of steps;
%   - compatibility and smoothness;
%   - existence and uniqueness;
%   - equilibria;
%   - scalar linear DDEs;
%   - transcendental characteristic equations;
%   - infinitely many characteristic roots;
%   - characteristic-root stability;
%   - delayed negative feedback;
%   - imaginary-axis crossings;
%   - critical-delay calculations;
%   - delay-induced oscillation;
%   - Hopf bifurcation;
%   - Lambert W representation;
%   - nonlinear DDE linearization;
%   - delayed logistic equations;
%   - feedback-control models;
%   - physiological and epidemiological delays;
%   - maturation delays;
%   - linear DDE systems;
%   - multiple-delay systems;
%   - neutral DDEs;
%   - distributed-delay characteristic equations;
%   - exponential memory kernels;
%   - state augmentation;
%   - relation to integral equations;
%   - Laplace transforms with history;
%   - numerical DDE methods;
%   - delayed-state interpolation;
%   - Euler approximation;
%   - discontinuity tracking;
%   - applications and exercises.
%
% Notes:
%   - Differential-Algebraic Equations is developed next in Section 9.8.
%   - Section 9.8 completes Chapter 9.
%   - Numerical DDE methods are developed more deeply in Chapter 11.
%
%%%%%%%%%%%%%%%%%%%%%%%%%%%%%%%%%%%%%%%%%%%%%%%%%%%%%%%%%%%%